\documentclass[letterpaper,12pt]{article}
\usepackage[T1]{fontenc}
\usepackage[utf8]{inputenc}
\usepackage[spanish]{babel}
\usepackage{framed}
\usepackage{amsfonts,setspace}
\usepackage{amsmath}
\usepackage{amssymb, amsmath, amsthm}
\usepackage{comment}
\usepackage[most]{tcolorbox}
\usepackage{amssymb}
\usepackage{dsfont}
\usepackage{float}
\usepackage{anysize}
\usepackage{pdfpages}
\usepackage{multicol}
\usepackage{enumerate}
\usepackage{graphicx}
\usepackage[dvipsnames]{xcolor}
\usepackage[left=1.5cm,top=0.9cm,right=1.5cm, bottom=1.4cm]{geometry}
\usepackage{fancyhdr}
\pagestyle{fancy}
\definecolor{blue(pigment)}{rgb}{0.2, 0.2, 0.6}
% Page
\oddsidemargin -0.4in
\textwidth 7in
\topmargin -0.6in
\textheight 9.1in
\parindent 0em
\parskip 1ex

%Teoremas, Lemas, etc.
\theoremstyle{plain}
\newtheorem{teo}{Teorema}
\newtheorem{lem}{Lemma}
\newtheorem{prop}{Proposici\'on}
\newtheorem{cor}{Corolario}
\theoremstyle{definition}
\newtheorem{defi}{Definici\'on}
\newtheorem{eje}{Ejemplo}
\newtheorem{obs}{Observaci\'on}
% fin macros

% Comandos más usados
% Conjuntos
\newcommand{\N}{\mathbb{N}}
\newcommand{\Z}{\mathbb{Z}}
\newcommand{\Q}{\mathbb{Q}}
\newcommand{\R}{\mathbb{R}}
\newcommand{\C}{\mathbb{C}}
\newcommand{\RR}{\mathbb{R} \times \mathbb{R}}
% Flechas
\newcommand{\ssi}{\Longleftrightarrow}
\newcommand{\imp}{\Longrightarrow}
\newcommand{\pmi}{\Longleftarrow}
% Trigonométricas
\newcommand{\sen}{\text{sen}}

\usepackage{versions}
\includeversion{solution}

\let\epsilon\varepsilon

\begin{document}
\def\Pauta{1}
%%%%%ESCRIBIR 1 si quieren ver la Prueba con Pauta%%%%
%%%%%ESCRIBIR 0 si quieren ver sólo la Ayudantía%%%%
	
	\fancyhead[L]{\itshape{Escuela de Ingeniería}}
	\fancyhead[R]{\itshape{Universidad de O'Higgins}}
	
	\begin{minipage}{11.5 cm}
         \vspace{-6mm}
		\begin{flushleft}
			\vspace{1mm}
			\textbf{ING1002-2 Cálculo Diferencial e Integral}\\
			\textbf{Profesor:}  Gonzalo Flores G.  \\
			\textbf{Ayudante:}  Juan I. Riquelme \\
		\end{flushleft}
	\end{minipage}
	
	\begin{picture}(2,3)
		\put(430,9.5){\includegraphics[scale=0.22]{Imagenes/UOH.png}}
	\end{picture}
	\vspace{-7mm}
	\begin{center}
		\LARGE \bf{Ayudantía 1: Límites}
	\end{center}
    \vspace{-5.5mm}
    \begin{center}
		28 de Agosto de 2025
	\end{center}

\begin{enumerate}[\textbf{P\arabic{enumi}.}]
\item \noindent Determine, si existe, el valor de los siguientes l\'{\i}mites:
\begin{enumerate}
\item $\displaystyle \lim_{x \to -1} \frac{x^3 + 1}{x^2 - 1}$

\vspace{0.2cm}
\color{blue}
    \noindent{\bf Sol:}
\[
\lim_{x \to -1} \frac{x^3 + 1}{x^2 - 1}
= \lim_{x \to -1} \frac{(x + 1)(x^2 - x + 1)}{(x + 1)(x - 1)}
= \lim_{x \to -1} \frac{x^2 - x + 1}{x - 1}
= \frac{(-1)^2 - (-1) + 1}{-1 - 1} = \frac{3}{-2} = -\frac{3}{2}
\]
\color{black}

\item $\displaystyle \lim_{x \to 0} x^3 \sen\left(\frac{1}{x}\right)$   

\vspace{0.2cm}
\color{blue}
    \noindent{\bf Sol:} 
\[
\text{Como } -1 \leq \sin\left(\frac{1}{x}\right) \leq 1, \text{ entonces } -x^3 \leq x^3 \sin\left(\frac{1}{x}\right) \leq x^3 
\]
Luego
\[
\lim_{x \to 0} -x^3 = 0 = \lim_{x \to 0} x^3 \implies
\lim_{x \to 0} x^3 \sin\left(\frac{1}{x}\right) = 0 \quad\quad \textcolor{blue}{\text{(Por teorema del Sandwich)}}
\]
\color{black}

\item $\displaystyle \lim_{x \to 0} \frac{\sqrt{x+2} - \sqrt{2}}{x}$

\vspace{0.2cm}
\color{blue}
\textbf{Sol:} Multiplicamos arriba y abajo por el conjugado del numerador $\sqrt{x+2} + \sqrt{2}$
\[
\Rightarrow \lim_{x \to 0} \frac{\sqrt{x+2} - \sqrt{2}}{x}
= \lim_{x \to 0} \frac{\sqrt{x+2} - \sqrt{2}}{x}  \cdot \left( \frac{\sqrt{x+2} + \sqrt{2}}{\sqrt{x+2} + \sqrt{2}} \right)
= \lim_{x \to 0} \frac{(x+2) - 2}{x(\sqrt{x+2} + \sqrt{2})}
\]
\[= \lim_{x \to 0} \frac{x}{x(\sqrt{x+2} + \sqrt{2})}
= \lim_{x \to 0} \frac{1}{\sqrt{x+2} + \sqrt{2}} = \frac{1}{\sqrt{2} + \sqrt{2}} = \frac{1}{2\sqrt{2}} \quad \quad \textcolor{blue}{(Evaluando)}
\]
\color{black}
\item $\displaystyle \lim_{x \to 0} \frac{1 - \cos(\pi x)}{3x^2} $

\vspace{0.2cm}
\color{blue}
\noindent{\bf Sol:} Usamos el límite notable
\[
\lim_{u \to 0} \frac{1 - \cos(u)}{u^2} = \frac{1}{2}
\]
Sea el cambio de variable \( u = \pi x \):
\begin{align*}
\lim_{x \to 0} \frac{1 - \cos(\pi x)}{3x^2}
&= \frac{1}{3} \cdot \lim_{x \to 0} \frac{1 - \cos(\pi x)}{x^2} = \frac{1}{3} \cdot \lim_{x \to 0} \left( \frac{1 - \cos(\pi x)}{(\pi x)^2} \cdot \pi^2 \right) \\
&= \frac{1}{3} \cdot \pi^2 \cdot \lim_{u \to 0} \frac{1 - \cos(u)}{u^2} = \frac{1}{3} \cdot \pi^2 \cdot \frac{1}{2} \\
&= \frac{\pi^2}{6}
\end{align*}
\color{black}
\end{enumerate}

\item \textbf{[Subida aparte]}
\item Considere la siguiente función $f$ definida en los reales sin el cero.
$$f(x) = \dfrac{x^2-1}{x^2+x}$$
    \begin{enumerate}
        \item Realice una pequeña tabla de valores de $f$ cerca del cero.
        
  Por ejemplo, evaluando \( f(x) \) en algunos valores positivos y negativos cercanos a cero (excluyendo \( x = 0 \), ya que no está en el dominio):

\color{blue} \textbf{Sol:} Evaluamos valores de $x$ cerca de 0 para conjeturar el límite.
        \[
        \begin{array}{|c|c|}
        \hline
        x & f(x) = \dfrac{x^2 - 1}{x^2 + x} \\
        \hline
        -0.1 & \dfrac{(-0.1)^2 - 1}{(-0.1)^2 + (-0.1)}  \approx 11 \\
        -0.01 & \dfrac{0.0001 - 1}{0.0001 - 0.01}   \approx 101 \\
        -0.001 & \dfrac{10^{-6} - 1}{10^{-6} - 0.001}   \approx 1001 \\
        \hline
        0.001 & \dfrac{10^{-6} - 1}{10^{-6} + 0.001}  \approx -999 \\
        0.01 & \dfrac{0.0001 - 1}{0.0001 + 0.01} \approx -99 \\
        0.1 & \dfrac{0.01 - 1}{0.01 + 0.1}\approx -9 \\
        \hline
        \end{array}
        \]
        Se observa que:
        \begin{multicols}{2}
            \begin{itemize}
            \item Cuando \( x \to 0^- \), \( f(x) \to +\infty \)
            \item Cuando \( x \to 0^+ \), \( f(x) \to -\infty \)
        \end{itemize}
        \end{multicols}
\color{black}
        
        \item Calcule las asintotas verticales y horizontales de la función $f$.

\color{blue}
\textbf{Sol:} 

        \begin{itemize}
        \item \textbf{Asíntotas verticales:}
        El denominador se anula cuando:
        \[
        x^2 + x = 0 \Rightarrow x(x + 1) = 0 \Rightarrow x = 0 \text{ o } x = -1
        \]
        Para determinar si hay asíntotas verticales en esos puntos, analizamos los límites laterales
        \begin{itemize}
            \item En \( x = 0 \):
            \[
            \lim_{x \to 0^-} f(x) = \lim_{x \to 0^-} \frac{x^2 - 1}{x^2 + x} = \frac{-1}{0^-} = +\infty
            \]
            \[
            \lim_{x \to 0^+} f(x) = \lim_{x \to 0^+} \frac{x^2 - 1}{x^2 + x} = \frac{-1}{0^+} = -\infty
            \]
            Por lo tanto, \( x = 0 \) es una \textbf{asíntota vertical}.
        \end{itemize}
        \item \textbf{Asíntotas Horizontales:}
        Una vez realizada la división de polinomios tenemos
        $$
        \frac{x^2 - 1}{x^2 + x} = 1- \frac{1}{x}
        $$
        Con 1 el cociente, luego $y=1$ es asíntota horizontal.
        \end{itemize}
\color{black}

        \item Con lo anterior y utilizando tabla de signos bosqueje $f$.

\color{blue}
\begin{center}
\begin{tabular}{|c|c|c|c|}
\hline
Intervalo & \( (-\infty, 0) \) & \( (0, 1) \) & \( (1, \infty) \) \\
\hline
\( x - 1 \) & \(-\) & \(-\) & \(+\) \\
\hline
\( x \)     & \(-\) & \(+\) & \(+\) \\
\hline
\(\dfrac{x - 1}{x} \) & \(+\) & \(-\) & \(+\) \\
\hline
\end{tabular}
\end{center}
A partir de la tabla de signos y que tenemos asíntota vertical en $x=0$ y asíntota horizontal en $y=1$, tenemos que la gráfica de la función es la siguiente
\begin{figure}[H]
    \centering
    \includegraphics[width=0.7\linewidth]{Imagenes/P3_Ayudantía_1_CDI.jpg}
    \caption{}
    \label{fig:placeholder}
\end{figure}

    \end{enumerate}


\item Considere la función $h:\R\to\R$ definida por

$$h(t) = \left\{
    \begin{array}{cc}
       ct^2+2t  & \text{Si $t<2$} \\
        t^3-ct & \text{Si $t\geq2$}
    \end{array}
\right.$$

donde $c$ es una constante a determinar.
    \begin{enumerate}
        \item Determine el valor de $c$ para que $\displaystyle \lim_{t \to 2} h(t)$ exista.

\color{blue}
    Para que el límite exista, debe cumplirse que:
    \[
    \lim_{t \to 2^-} h(t) = \lim_{t \to 2^+} h(t)
    \]
    Calculamos los límites laterales:
    \[
    \lim_{t \to 2^-} h(t) = \lim_{t \to 2^-} (ct^2 + 2t) = 4c + 4
    \]
    \[
    \lim_{t \to 2^+} h(t) = \lim_{t \to 2^+} (t^3 - ct) = 8 - 2c
    \]
    Igualando ambos:
    \[
    4c + 4 = 8 - 2c \Rightarrow 6c = 4 \Rightarrow c = \frac{2}{3}
    \]

\color{black}

        \item Bosqueje $h$.

\color{blue}
\textbf{Sol:} Una vez tenemos que $c=\frac{2}{3}$, gráficamos $h(t)$ para $t<2$ y $t>2$.
\color{black}

\begin{figure}[H]
    \centering
    \includegraphics[width=0.7\linewidth]{Imagenes/P4_Ayudantía_1_CDI.jpg}
    \caption{}
    \label{fig:placeholder}
\end{figure}

        
    \end{enumerate}

\end{enumerate}


\end{document}

Por la identidad fundamental: $\sen^2(x)+\cos^2(x)=1$
$$ \lim_{x\to0} \frac{1 - \cos^{2}(\pi x)}{3x^{2}} = \lim_{x\to0} \frac{\sin^{2}(\pi x)}{3x^{2}}= \lim_{x\to0} \frac{\sin^{2}(\pi x)}{3x^{2}} \cdot \textcolor{red}{\frac{\pi}{\pi}}= \lim_{x\to0} \frac{\sin^{2}(\pi x)}{3x^{2}}=\frac{1}{3}\lim_{u\to0}\left(\frac{\sin(\pi x)}{\pi x}\right)^2\pi $$
Usando el l\'imite notable $\displaystyle \lim_{x \to 0}\frac{\sen(x)}{x}=1$:
$$
=\frac{1}{3}\left(\lim_{u\to0}\frac{\sin(\pi x)}{\pi x}\right)^2\pi =\frac{1}{3} \pi
$$
\color{black}

\item $\displaystyle 
\lim_{x \to 3} \frac{9 - x^2}{5 - \sqrt{x^2 + 16}}
$

\vspace{0.2cm}
\color{blue}
\textbf{Sol:}
\begin{align*}
\lim_{x \to 3} \frac{9 - x^2}{5 - \sqrt{x^2 + 16}} 
&= \lim_{x \to 3} \frac{(3 - x)(3 + x)}{5 - \sqrt{x^2 + 16}} = \lim_{x \to 3} \frac{-(x - 3)(x + 3)}{5 - \sqrt{x^2 + 16}} \\
&\text{Multiplicamos por el conjugado: } \frac{5 + \sqrt{x^2 + 16}}{5 + \sqrt{x^2 + 16}} \\
&= \lim_{x \to 3} \frac{-(x - 3)(x + 3)(5 + \sqrt{x^2 + 16})}{(5 - \sqrt{x^2 + 16})(5 + \sqrt{x^2 + 16})} = \lim_{x \to 3} \frac{-(x - 3)(x + 3)(5 + \sqrt{x^2 + 16})}{25 - (x^2 + 16)} \\
&= \lim_{x \to 3} \frac{-(x - 3)(x + 3)(5 + \sqrt{x^2 + 16})}{9 - x^2} = \lim_{x \to 3} \frac{-(x - 3)(x + 3)(5 + \sqrt{x^2 + 16})}{-(x^2 - 9)} \\
&= \lim_{x \to 3} (5 + \sqrt{x^2 + 16}) = 5 + \sqrt{9 + 16} \\&= 5 + \sqrt{25} = 5 + 5 = 10
\end{align*}
\color{black}

\item $\displaystyle \lim_{x \to +\infty} \left( \sqrt{3x^2 + 8x + 6} - \sqrt{3x^2 + 3x + 6} \right)$   

 \vspace{0.2cm}
 \color{blue}
     \noindent{\bf Sol:}
 \[
 \lim_{x \to +\infty} \left( \sqrt{3x^2 + 8x + 6} - \sqrt{3x^2 + 3x + 6} \right)
 = \lim_{x \to +\infty} \frac{5x}{\sqrt{3x^2 + 8x + 6} + \sqrt{3x^2 + 3x + 6}}
 \]
 \[
 = \lim_{x \to +\infty} \frac{5}{\sqrt{3 + \frac{8}{x} + \frac{6}{x^2}} + \sqrt{3 + \frac{3}{x} + \frac{6}{x^2}}}
= \frac{5}{2\sqrt{3}}
 \]
\color{black}

    \item $\displaystyle \lim_{x \to +\infty} \frac{e^{3x} - e^{-3x}}{e^{3x} + e^{-3x}}$
    
 \vspace{0.2cm}
 \color{blue}
  \noindent{\bf Sol:} Vamos a multiplicar el númerador y el denominador por $e^{-3x}$
 $$
\lim_{x \to +\infty} \frac{e^{3x} - e^{-3x}}{e^{3x} + e^{-3x}} \cdot \left(\dfrac{e^{-3x}}{e^{-3x}}\right) = \lim_{x \to +\infty} \frac{1 - e^{-6x}}{1 + e^{-6x}}
$$
Evaluando
\[
 = \frac{1 - 0}{1 + 0} = 1
\]
\color{black}